\documentclass[10pt,a4paper]{amsart}
\usepackage[margin=19mm]{geometry}
\usepackage{amsmath,amssymb,mathtools}
\usepackage[scr=rsfs,cal=euler]{mathalfa}
\usepackage{xfrac}
\usepackage[unicode,hidelinks,bookmarks=false]{hyperref}
\newcommand{\ie}{i.e.\ }
\newcommand{\cf}{cf.\ }
\newcommand{\resp}{respectively\ }
\newcommand{\Subsection}{\subsection}
\providecommand{\nobreakdash}{\nobreak\hbox{-}}
\setlength{\emergencystretch}{2em}
\multlinegap0pt
\usepackage[shortlabels]{enumitem}
\usepackage{amssymb}
\usepackage{xcolor}
\usepackage{tikz}
\newtheorem{theorem}{Theorem}
\renewcommand{\ge}{\geqslant}
\renewcommand{\le}{\leqslant}
\newcommand{\norm}[1]{\lVert #1 \rVert}
\title{Primariness and the Primary Factorisation Property}
\author{Antonio Acuaviva, Tomasz Kania}
\date{Source: arXiv:2605.21711v1 (CC BY 4.0)}
\begin{document}
\maketitle

\noindent\emph{Source and licence.} Primariness and the Primary Factorisation Property. Version arXiv:2605.21711v1 (CC BY 4.0), distributed by arXiv under the Creative Commons Attribution 4.0 International licence, http://creativecommons.org/licenses/by/4.0/, as recorded in the arXiv licence metadata for this exact version and shown on the abstract page https://arxiv.org/abs/2605.21711v1. Full licence notice: \texttt{licenses/arxiv-2605.21711-CC-BY-4.0.txt}.\par\smallskip

\noindent\textit{Editorial scope.} Counted unit: the article's theorem thm:quant-pelc (source0.tex lines 4257--4283). The article's complete original proof of it is delivered with it (source0.tex lines 4285--4494, one \texttt{proof} environment of 210 lines). No mathematics is written, repaired or re-derived: the statement and the proof are the article's own lines, in the article's own notation, and no retained line is substituted or re-flowed.  No further source-local context is needed: every cross-reference of every retained line resolves inside the two delivered spans.  Nothing was omitted from any retained span, so the package carries no omission record.  No retained line contains a citation, so the wrapper carries no bibliography and no bibliography entry.  No retained line contains a figure environment, an image-inclusion command or drawing code, so no image is delivered and the package declares no asset dependency.  Editorial formatting only, in addition: an AMS \texttt{amsart} preamble that restates the article's own declarations and macros used by the retained lines, namely the environments \texttt{theorem} on separate counters, together with the standard packages those lines need.  The retained environments print in this document as Theorem 1 in order of appearance, whereas the article prints the same items with the numbering of its own full document.  The complete native source of the article as distributed in the arXiv e-print for this exact version is bundled unchanged as \texttt{sources/201078/source0.tex}.
\begin{theorem}[A quantitative form of Pe{\l}czy\'nski's decomposition method]\label{thm:quant-pelc}
Let $X$ and $Y$ be Banach spaces such that
\begin{enumerate}[label=\textup{(\roman*)}]
    \item $X$ has a $C_1$-complemented subspace $E$ which is $C_2$-isomorphic to $Y$;
    \item $Y$ has a $C_3$-complemented subspace $G$ which is $C_4$-isomorphic to $X$.
\end{enumerate}
Suppose further that either
\begin{enumerate}[label=(\alph*), ref=(\alph*)]
    \item \label{item:quant-pelc-a}
    there exists $1 \le r \le \infty$ such that
    \[
        X \simeq_{C_5} X \oplus_r X
        \qquad\text{and}\qquad
        Y \simeq_{C_6} Y \oplus_r Y,
    \]
    or
    \item \label{item:quant-pelc-b}
    \[
        X \simeq_{C_5} c_0(X)
        \qquad\text{or}\qquad
        X \simeq_{C_5} \ell_p(X)
    \]
    for some $1 \le p \le \infty$.
\end{enumerate}
Then there exists a constant $C \ge 1$, depending only on $C_1,\dots,C_6$ in case \ref{item:quant-pelc-a}, and only on $C_1,\dots,C_5$ in case \ref{item:quant-pelc-b}, such that $X \simeq_C Y$.
Moreover, the proof below gives an explicit admissible choice of $C$.
\end{theorem}

\begin{proof}
Choose a projection $P \colon X \to E$ with $\norm{P} \le C_1$, a projection $Q \colon Y \to G$ with $\norm{Q} \le C_3$, an isomorphism $S \colon Y \to E$ with $\norm{S}\norm{S^{-1}} \le C_2$, and an isomorphism $T \colon X \to G$ with $\norm{T}\norm{T^{-1}} \le C_4$.
After rescaling $S$ and $T$, we may assume that
\[
    \norm{S}=\norm{S^{-1}} \le \sqrt{C_2},
    \qquad
    \norm{T}=\norm{T^{-1}} \le \sqrt{C_4}.
\]
Set
\[
    F = \ker P,
    \qquad
    H = \ker Q.
\]

We first show that, for each fixed $1 \le r \le \infty$, there exist constants $a=a(C_1,C_2)$ and $b=b(C_3,C_4)$ such that
\(
    X \simeq_a Y \oplus_r F,
\)
and
\(
    Y \simeq_b X \oplus_r H.
\)
To this end, define
\[
    \Phi_X \colon X \to Y \oplus_r F,
    \qquad
    \Phi_X x = \bigl(S^{-1}Px,(I_X-P)x\bigr),
\]
and
\[
    \Psi_X \colon Y \oplus_r F \to X,
    \qquad
    \Psi_X(y,f)=Sy+f.
\]
Since $f \in F=\ker P$, we have $Pf=0$, and hence
\[
    \Psi_X\Phi_X x = S(S^{-1}Px)+(I_X-P)x = Px+(I_X-P)x = x.
\]
Also, if $y \in Y$ and $f \in F$, then
\[
    \Phi_X\Psi_X(y,f)=\bigl(S^{-1}P(Sy+f),(I_X-P)(Sy+f)\bigr)=(y,f),
\]
because $P(Sy)=Sy$ and $Pf=0$.
Thus $\Phi_X$ and $\Psi_X$ are inverse isomorphisms. Moreover,
\[
    \norm{\Phi_X x}_r \le \norm{S^{-1}Px}+\norm{(I_X-P)x}
    \le \bigl(\sqrt{C_2}\,C_1+1+C_1\bigr)\norm{x}.
\]
Likewise,
\[
    \norm{\Psi_X(y,f)} \le \norm{Sy}+\norm{f}
    \le (\sqrt{C_2}+1)(\norm{y}+\norm{f})
    \le 2(\sqrt{C_2}+1)\norm{(y,f)}_r.
\]
Hence
\(
    X \simeq_a Y \oplus_r F,
\)
where
\(
    a := 2(\sqrt{C_2}+1)\bigl(\sqrt{C_2}\,C_1+1+C_1\bigr).
\)

Similarly, defining the operators 
\(
    \Phi_Y \colon Y \to X \oplus_r H, 
    \Phi_Y y = \bigl(T^{-1}Qy,(I_Y-Q)y\bigr),
\)
and
\(
    \Psi_Y \colon X \oplus_r H \to Y, 
    \Psi_Y(x,h)=Tx+h,
\)
one obtains
\(
    Y \simeq_b X \oplus_r H,
\)
where
\(
    b := 2(\sqrt{C_4}+1)\bigl(\sqrt{C_4}\,C_3+1+C_3\bigr).
\)

We now treat the two cases separately.

\smallskip
\noindent
\textit{Case \ref{item:quant-pelc-a}.}
Fix $r$ as in the hypothesis.
Since $X \simeq_a Y \oplus_r F$, we obtain
\[
    X \simeq_a Y \oplus_r F
    \simeq_{C_6} (Y \oplus_r Y)\oplus_r F
    \simeq_1 Y \oplus_r (Y \oplus_r F)
    \simeq_a Y \oplus_r X,
\]
and therefore
\(
    X \simeq_{a^2C_6} Y \oplus_r X.
\)
Likewise,
\[
    Y \simeq_b X \oplus_r H
    \simeq_{C_5} (X \oplus_r X)\oplus_r H
    \simeq_1 X \oplus_r (X \oplus_r H)
    \simeq_b X \oplus_r Y,
\]
so
\(
    Y \simeq_{b^2C_5} X \oplus_r Y.
\)
Combining the last two displays, and using the canonical isometry $Y \oplus_r X \simeq_1 X \oplus_r Y$, we get
\[
    X \simeq_{a^2C_6} Y \oplus_r X
    \simeq_1 X \oplus_r Y
    \simeq_{b^2C_5} Y.
\]
Thus
\(
    X \simeq_{a^2b^2C_5C_6} Y.
\)

\smallskip
\noindent
\textit{Case \ref{item:quant-pelc-b}.}
We treat the case
\[
    X \simeq_{C_5} \ell_p(X),
    \qquad 1 \le p \le \infty;
\]
the case $X \simeq_{C_5} c_0(X)$ is identical, replacing $\ell_p$ by $c_0$ and $\oplus_p$ by $\oplus_\infty$ throughout.

Since $\ell_p(X)$ is canonically isometric to $X \oplus_p \ell_p(X)$, we have
\[
    X \simeq_{C_5} \ell_p(X)
    \simeq_1 X \oplus_p \ell_p(X)
    \simeq_{C_5} X \oplus_p X,
\]
and so
\(
    X \simeq_{C_5^2} X \oplus_p X.
\)

Now take $r=p$ in the first part of the proof.
Then
\[
    X \simeq_a Y \oplus_p F,
    \qquad
    Y \simeq_b X \oplus_p H.
\]
We next show that
\(
    Y \simeq_{(abC_5)^2} Y \oplus_p Y.
\)
Starting from $Y \simeq_b X \oplus_p H$, we obtain
\[
    Y \simeq_b X \oplus_p H
    \simeq_{C_5} \ell_p(X)\oplus_p H
    \simeq_a \ell_p(Y \oplus_p F)\oplus_p H
    \simeq_1 \ell_p(Y)\oplus_p \ell_p(F)\oplus_p H.
\]
Since $\ell_p(Y)$ is canonically isometric to $Y \oplus_p \ell_p(Y)$, it follows that
\[
    Y \simeq_{abC_5} Y \oplus_p \bigl(\ell_p(Y)\oplus_p \ell_p(F)\oplus_p H\bigr).
\]
On the other hand,
\[
    \ell_p(Y)\oplus_p \ell_p(F)\oplus_p H
    \simeq_1 \ell_p(Y \oplus_p F)\oplus_p H
    \simeq_a \ell_p(X)\oplus_p H
    \simeq_{C_5} X \oplus_p H
    \simeq_b Y.
\]
Therefore
\(
    Y \simeq_{(abC_5)^2} Y \oplus_p Y.
\)

We now repeat the same decomposition argument with $\oplus_p$ in place of $\oplus_r$.
Using $X \simeq_a Y \oplus_p F$, we get
\[
    X \simeq_a Y \oplus_p F
    \simeq_{(abC_5)^2} (Y \oplus_p Y)\oplus_p F
    \simeq_1 Y \oplus_p (Y \oplus_p F)
    \simeq_a Y \oplus_p X,
\]
whence
\(
    X \simeq_{a^2(abC_5)^2} Y \oplus_p X.
\)
Similarly, since $X \simeq_{C_5^2} X \oplus_p X$, we have
\[
    Y \simeq_b X \oplus_p H
    \simeq_{C_5^2} (X \oplus_p X)\oplus_p H
    \simeq_1 X \oplus_p (X \oplus_p H)
    \simeq_b X \oplus_p Y,
\]
and therefore
\(
    Y \simeq_{b^2C_5^2} X \oplus_p Y.
\)
Combining the last two displays, and using the canonical isometry $Y \oplus_p X \simeq_1 X \oplus_p Y$, we obtain
\(
    X \simeq_{a^2(abC_5)^2} Y \oplus_p X
    \simeq_1 X \oplus_p Y
    \simeq_{b^2C_5^2} Y.
\)
Thus $X \simeq_{a^4b^4C_5^4} Y$.
This completes the proof.
\end{proof}

\end{document}
